% MEASUREMENT_MAP.tex, a companion to PROOFS.tex (draft of 31 August 2026,
% as committed 2 September 2026). Written 6 October 2026, revised the same day
% to follow notes/writing_discipline.md.
%
% PURPOSE
%   A reference note, not part of the paper. For every input and every output
%   of the model in PROOFS.tex it records what the object is, what it
%   represents, whose it is, and how an observer could obtain it. It makes no
%   arguments and claims nothing about identification.
%
% PASTING ROWS INTO PROOFS.tex
%   Rows use only packages already loaded in the PROOFS.tex preamble (amsmath,
%   booktabs, array, longtable, natbib) and no macros of their own. Results are
%   named in words. On pasting, replace each name on the left with the
%   reference on the right.
%     "the payoff"                         -> \eqref{eq:payoff}
%     "burden-monotonicity"                -> \eqref{eq:burdenmono}
%     "the entry condition"                -> \eqref{eq:entry}
%     "Proposition (anonymity)"            -> Proposition~\ref{prop:anon}
%     "Proposition (count invariance)"     -> Proposition~\ref{prop:countinv}
%     "Proposition (when identities are pinned)" -> Proposition~\ref{prop:idunique}
%     "Proposition (assortative selection)" -> Proposition~\ref{prop:assort}
%     "Corollary (monotonicity)"           -> Corollary~\ref{cor:monotone}
%     "P1", "P2", "P6", "P7", "P8", "P-MU" -> \S\ref{prop:P1}, \S\ref{prop:P2},
%                                             \S\ref{prop:P6}, \S\ref{prop:P7},
%                                             \S\ref{prop:P8}, \S\ref{prop:PMU}
%     "Theorem (key result)" (P3)          -> \S\ref{prop:P3}
%     "Theorem (P9)"                       -> Theorem~\ref{thm:P9}
%     "Theorem (P9, general form)"         -> Theorem~\ref{thm:P9gen}
%     "Proposition (endogenous margin)"    -> Proposition~\ref{prop:endogenous}
%     "Proposition (margin condition)"     -> Proposition~\ref{prop:tailprop}
%     "the band"                           -> \S\ref{prop:band}
%     "Proposition (strictness)"           -> Proposition~\ref{prop:strict}
%     "the R1 subsection", "the R2 subsection" -> \S\ref{sec:R1}, \S\ref{sec:R2}
%     "the scope paragraph"                -> \S\ref{sec:primitives}
%   The two citations need Schoemaker1982 and Sen1973 added to refs.bib.
%
% VERIFICATION NOTE
%   Each row that asserts mathematics is listed with the place it was checked.
%   "Manuscript" means the result is restated from PROOFS.tex and carries that
%   document's own status and tier (its Table "Verification index"). "MMn" means
%   the Lean theorems of that name in lean/mathlib/MeasurementMap.lean (pass 3h,
%   7 Oct 2026), audited by checks/verify_mathlib.py. The earlier SymPy version
%   remains in checks/verify_measurement_map.py and no longer counts as
%   evidence. Manuscript rows marked (S) below have since moved to Lean: P1
%   (RepresentationFOSD.lean), P3 (StepIdentity.lean), P6 (SaturationBenchmark.lean,
%   p6_limit), the divergence of kappa (KappaSpread.lean), P-MU (PMU.lean).
%     Framing and the u, kappa, V rows. The cardinal reading is the manuscript's
%       own ("cardinal utility index", N4 discussion). The entry condition is the
%       expected-payoff comparison (MM9). The entry condition is invariant under
%       (u,V)->(au+b,aV) (MM6a-c), and its invariants are exactly the functions
%       of kappa/V (MM6d). An expected-utility ranking is invariant under aU+b
%       and not under a monotone non-affine map (MM10).
%     Inputs, V. Scale factor of Delta, manuscript P1 (S).
%     Inputs, mu. sup supp(G)=mu needs r to have upper endpoint 1 (MM8).
%     Inputs, kappa. Strictly decreasing and divergent at c are manuscript
%       assumptions (burden-monotonicity; S16).
%     Inputs, F and G. F<=G is manuscript P2 (L).
%     Inputs, x0 and lambda. Recovery formulas (MM1, MM2).
%     Outputs, k*. Count invariance, manuscript (L).
%     Outputs, w_(k*). The rank-k* challenger can be absent from an equilibrium
%       (MM7, on the manuscript's own example).
%     Outputs, Delta. P1 representation, manuscript (S).
%     Outputs, step. P3, manuscript (S).
%     Outputs, P6. With the factor V (MM3); manuscript S9 without it.
%     Outputs, P7. With the factor V (MM4); manuscript S18-S20 and
%       entry_index_bounded.
%     Outputs, P8, P9, P9-gen, endogenous margin, margin condition. Manuscript (L).
%     Outputs, band. Manuscript (L).
%     Outputs, P-MU. With the factor V (MM5); manuscript S10 without it.
%     Outputs, k*c. Identical across equilibria, manuscript N4 discussion.
%     Outputs, U_i and the two sums. Dependence on a and b, and the sign of the
%       difference between two equilibria preserved (MM6e-f).
%     Recovery 1-2 (MM1, MM2). Recovery 4 restates the equilibrium conditions of
%       Proposition (count invariance), divided by V>0.
%
% SOURCES
%   Schoemaker (1982), Journal of Economic Literature 20(2), pp. 529-563, read
%   in full [F] on 6 Oct 2026 from Drive folder 1r7J3xJ2rotpmlRaBANCNhviy2d8a1bx_,
%   file 1982-Schoemaker_ExpectedUtilityModel-VariantsPurposesLimits_JofEconomicLiterature.pdf.
%   Quotations with pages.
%     "unique up to positive linear transformations ... U*(x)=aU(x)+b for
%     numbers a>0 and b" (p.531)
%     "utility, in the NM context, is used to represent preferences whereas in
%     neoclassical theory it determines (or precedes) preference" (p.532)
%     "from a preference perspective, NM utility theory is ordinal in that it
%     provides no more than ordinal rankings of lotteries" (p.533)
%     "NM utility should not be interpreted as measuring strength of
%     preference under certainty" (p.533)
%     "a concave u(x) might erroneously be interpreted as implying that equal
%     increments in money (under certainty) contribute to utility at a
%     decreasing rate" (p.535)
%   That this model is of his NM kind rests on MM9, the entry condition being
%   the expected-payoff comparison, and not on anything the manuscript says.
%   MM10 instantiates his p.531 claim and its failure under a monotone
%   non-affine map. The phrase "not quantities in the world" is this note's
%   stance, and Schoemaker supports it only as far as p.532 goes.
%   Sen (1973), Economica, New Series, 40(159), pp. 241-259, read in full [F]
%   on 6 Oct 2026 from the same Drive folder, file "sen73-behaviour and the
%   concept of preference.pdf" (JSTOR copy). Quotations with pages.
%     "the underlying relation in terms of which individual choices can be
%     explained" (p.253)
%     "will simply be the binary representation of individual choice" (p.253)
%     "a numerical representation of the as if preference cannot be
%     interpreted as individual welfare" (p.254)
%     "information need not be restricted to distant observations of choices
%     made" (p.257)
%     "well-known limitations of the questionnaire method" (p.258)
%     "it is not in general possible to guarantee both simultaneously" (p.259)
%   Sen (1971), Review of Economic Studies 38(3), pp. 307-317, is not cited,
%   because no full text was accessible (the Scribd preview shows metadata
%   only and OUP is paywalled). Sen (1973, p.253 fn.) cites it for the
%   conditions for binariness and transitivity of revealed preference, and that
%   content is not used here.
%
% DEPENDENCE ON UNAPPLIED CHANGES
%   No row depends on an unapplied change. Two rows will be needed later. The
%   welfare section (TODO.md W1) is unwritten, so no welfare objective appears
%   in the Outputs table. If the incumbent's entry is endogenised (TODO.md M1),
%   C moves from Inputs to Outputs.
%   PROOFS.tex writes P6, the P7 proof and P-MU without the factor V. The rows
%   below carry V explicitly and need no change if the manuscript adds it.

\documentclass[11pt]{article}
\usepackage[T1]{fontenc}
\usepackage[utf8]{inputenc}
\usepackage{mathptmx}
\usepackage{microtype}
\usepackage[a4paper,margin=2cm]{geometry}
\usepackage{amsmath,amssymb}
\usepackage{booktabs}
\usepackage{array}
\usepackage{longtable}
\usepackage{natbib}

\newcolumntype{P}[1]{>{\raggedright\arraybackslash}p{#1}}

\title{Threshold Entry in a Fixed-Prize Contest: Measurement Map}
\author{Companion note to the proofs document}
\date{6 October 2026}

\begin{document}
\maketitle

\noindent The model is a static, one-shot contest for one prize of value $V$
among $Q$ wealth-heterogeneous challengers and one incumbent. Each challenger
chooses $e_i\in\{0,1\}$, paying the nominal cost $c$ to invest; scores
$\sigma_i=\mu r_i+(1-\mu)s_ie_i$ are then drawn, and the highest score takes the
prize. The realised wealth profile is common knowledge, so the model contains
no beliefs, priors or posteriors, and no observer-side parameter such as
measurement precision. The discrete objects are $Q$, the incumbent's label
$C\in\{F,G\}$, the decisions $e_i$, the entrant set $S$, the count $k^*$, the
index $m$ and the rank $j$. The continuous objects are wealth, $c$, $V$, $\mu$,
the draws, the scores and the functions $u$, $\kappa$, $F$, $G$, $\varphi$,
$H_m$ and $\Delta$; the wealth law $\Lambda$ and the laws of $r$ and $s$ enter
as distributions. The utility of wealth $u$, the prize value $V$, and every
object built from them ($\kappa$, $\Delta$ in its own units, $U_i$, and the
sums of these across agents) are the modeller's constructs, attributed to the
agents to rationalise recorded choices, rather than quantities in the world.
The entry decision is a choice between lotteries, because the entry condition
is exactly the comparison of the expected payoffs of investing and of not
investing. The model is therefore of the von Neumann--Morgenstern kind in the
sense of \citet{Schoemaker1982}. In a model of that kind utility ``is used to
represent preferences'' rather than determining them (p.~532), and it is
unique up to $aU+b$ with $a>0$ (p.~531). It is cardinal in the measurement
sense only, since as a preference theory it ``provides no more than ordinal
rankings of lotteries'' and ``should not be interpreted as measuring strength
of preference under certainty'' (p.~533). \citet{Sen1973} draws the matching
line on the side of choice. Preference may be defined ``as the underlying
relation in terms of which individual choices can be explained'', and it then
``will simply be the binary representation of individual choice'' (p.~253).
Reading the same relation also as welfare is a further step, and ``it is not
in general possible to guarantee both simultaneously'' (p.~259). Accordingly
the entry condition is unchanged under $(u,V)\mapsto(au+b,\,aV)$ with $a>0$,
and a statement about these constructs means something only if it survives
that transformation. The last column separates five kinds of object, namely
what an observer records, what is computed from records, facts of nature that
exist but are never observed, constructs together with the recorded choices
they rationalise, and computed references. Numerical example parameters (the
CRRA coefficient, the witnesses, the constructions in the claim on
burden-monotonicity) are not inputs and are omitted, as are the refuted claims
R1 and R2 except where the manuscript reuses their objects.

\section*{Inputs}

\footnotesize
\begin{longtable}{@{}P{2.0cm}P{3.1cm}P{5.1cm}P{5.5cm}@{}}
\caption{Inputs of the model.}\\
\toprule
\textbf{Symbol} & \textbf{Type} & \textbf{Represents} & \textbf{Obtained how} \\
\midrule
\endfirsthead
\toprule
\textbf{Symbol} & \textbf{Type} & \textbf{Represents} & \textbf{Obtained how} \\
\midrule
\endhead
\bottomrule
\endfoot
$Q$ & Positive integer. & The number of challengers, ``$Q$ challengers plus
  one incumbent''. A fact of the contest. & Recorded by the observer as the
  number of challengers in the contest. \\ \addlinespace
$C\in\{F,G\}$ & One of two CDFs; a binary label. & The incumbent's score CDF,
  $F$ if she invests and $G$ if she does not. ``The additional player,
  entering exogenously; appears only through her score CDF $C$ in $H_m$.'' &
  Recorded by the observer as whether the incumbent paid $c$. A fixed input of
  the model rather than a choice made in it. \\ \addlinespace
$w_i$; $(w_1,\dots,w_Q)$; $w_{(j)}$, $j$ & Real number, $w_i>c$; a vector in
  $\mathbb R^Q$; its descending order statistics; integer rank, $1$ richest. &
  ``Challenger $i$'s resource.'' The realised profile is common knowledge
  among the challengers. A fact about each challenger, known to all of them. &
  Recorded by the observer before the entry decisions. Order statistics and
  ranks are computed from the record. The record is the observer's, whereas
  the challengers' knowledge of the profile is the manuscript's assumption and
  is not shown by the record. \\ \addlinespace
$\Lambda$ & A distribution across challengers, support bounded below by
  $c$. & ``The law from which the profile is drawn.'' A fact of nature. &
  Never observed as a law; each recorded profile is a draw from it. It enters
  the results only through the realised order statistics, the mean
  $\bar\mu$, the value $1-\Lambda(\bar\mu)$ and quantile functions. \\
  \addlinespace
$c$ & Real number, $c>0$, identical for all agents. & ``The nominal cost,
  identical for all agents'', paid ``to acquire social capital''. In
  resource units, and what society gives up. & Recorded by the observer as the
  outlay required to invest. \\ \addlinespace
$u$ & Function of wealth, increasing, satisfying burden-monotonicity; a
  cardinal index, fixed only up to $u\mapsto au+b$, $a>0$, together with
  $V\mapsto aV$. & ``The utility of wealth.'' The modeller's construct,
  attributed to the agent, rather than a quantity in the world. & Not
  observed, being a construct. It rationalises the recorded entry decisions,
  and enters the model only through $\kappa$. \\ \addlinespace
$\kappa(w)=u(w)-u(w-c)$ & Function on $(c,\infty)$; continuous, strictly
  decreasing, $\kappa\to\infty$ as $w\downarrow c$; scales with $a$ under the
  transformation above. & ``The burden of $c$ in utility units'', the entry
  cost (utility). ``The sole route by which the wealth distribution affects
  behaviour.'' The modeller's construct, attributed to the agent. It is a
  difference of a utility that represents preferences over lotteries, rather
  than a measure of how strongly the cost is felt (\citealp{Schoemaker1982},
  pp.~533,~535). & Not observed, being a construct. It rationalises the
  recorded entry decisions. Only the ratio $\kappa/V$ survives the
  transformation, and recorded entry decisions restrict that ratio to a set
  rather than to a value (Recovery~4). \\ \addlinespace
$V$ & Real number, $V>0$, in the units of $u$; scales with $a$ under the
  transformation above. & The prize's value. The prize is ``exogenous and
  fixed'', rank-allocated, a non-market good, additively separable and
  ``valued identically by rich and poor''. The modeller's construct,
  attributed to the agents. & Not observed, being a construct. It enters only
  as a scale factor of $\Delta$ (P1). The prize itself is recorded. Where two
  prizes are compared (P8), the recorded datum is the agents' choice between
  them, or their answer to a question ranking them, made before entry. An
  ordering of $V$ is what rationalises either record. \citet{Sen1973},
  pp.~257--258, counts answers to questions as information alongside
  behaviour, noting ``well-known limitations of the questionnaire
  method''. \\ \addlinespace
$\mu$ & Real number; the manuscript works with $\mu\in(0,1)$ and the limit
  $\mu\to0$. & The weight of $r$ in the score. A fact of the contest. & Never
  observed; enters only through $F$ and $G$. The R1 subsection states
  $\sup\operatorname{supp}(G)=\mu$, which holds when the support of $r$ has
  upper endpoint $1$, a condition the primitives do not state. \\
  \addlinespace
$r_i$ & Real-valued draw; independent across players and of wealth;
  arbitrary distribution. & The score component every player draws, whether
  or not she invests. A fact of nature. & Never observed separately. For a
  non-investor the recorded score is $\mu r_i$. \\ \addlinespace
$s_i$ & Real-valued draw, $s_i\ge0$; independent across players and of
  wealth; arbitrary distribution. & The component that investing adds, which
  the scope paragraph calls ability. A fact of nature. & Never observed
  separately. For an investor the recorded score is
  $\mu r_i+(1-\mu)s_i$. \\ \addlinespace
$F$, $G$ & CDFs; continuous, common lower endpoint; $F\le G$ pointwise
  (P2). & ``Law of $\sigma$ for an entrant ($F$) and an outsider ($G$).'' Facts
  of nature. ``All results are stated in $F$ and $G$ only.'' & Computed from
  records as the distributions of recorded scores among investors and among
  non-investors, pooled across challengers. Pooling across wealth uses the
  independence of the draws from wealth (Proposition (anonymity)). \\
  \addlinespace
$\varphi=G-F$ & Function, $\varphi\ge0$, vanishing at both endpoints of the
  union of supports. & ``The wedge investment buys.'' & Computed only, from
  $F$ and $G$. \\ \addlinespace
$m$ & Integer in $\{0,\dots,Q-1\}$. & The number of other challengers who
  invest, as faced by a challenger deciding. The rank-$j$ challenger,
  entering as the $j$-th entrant, faces $m=j-1$. & Computed from the rank. \\
  \addlinespace
$\theta$ & Real number; $\theta=0$ is the model. & Tilt of the law of $s$
  against wealth rank, outside the model's scope, used in the scope paragraph
  to show where count invariance fails. & Never observed as a parameter. The
  independence that $\theta$ relaxes is a direct comparison of recorded
  investor scores across wealth ranks. \\ \addlinespace
$w\mapsto w'$; $T$ & A map between wealth profiles; $T$ increasing for a
  pivot-spread. & ``Any map between wealth profiles'', the change in the
  wealth distribution that a comparative static is about. & Recorded as two
  wealth profiles, before and after a change, or of two contests. $T$ is seen
  only at the recorded wealths. \\ \addlinespace
$x_0$ & Real number, in wealth units. & The pivot, ``the wealth at which the
  displacement changes sign.'' & Computed from two recorded profiles. In
  general it is located only between two adjacent ranks; in the linear family
  it is given by Recovery~1. \\ \addlinespace
$\lambda$ & Real number, $\lambda>1$ (Theorem (P9)), $\lambda\ge1$ where the
  manuscript checks the pivot-spread property. & The factor of the linear spread
  $T_\lambda(w)=x_0+\lambda(w-x_0)$. & Computed from two recorded profiles
  (Recoveries 1--2). \\ \addlinespace
$\bar\mu$ & Real number, $\bar\mu=\mathbb E[w]$. & Mean wealth, and the pivot
  of the linear mean-preserving spread of Theorem (P9). & Computed from
  recorded wealths (Recovery~2). \\ \addlinespace
$\Lambda_\alpha$, $\Lambda_\beta$; $r_j$ & Two wealth distributions; quantile
  levels. & The quantile coupling, which reads the displacement as
  $\Lambda_\beta^{-1}(r_j)-\Lambda_\alpha^{-1}(r_j)$. & The laws are never
  observed; the displacement is computed from two recorded sorted profiles.
  Here $r_j$ is a quantile level rather than the draw $r_i$. \\
\end{longtable}
\normalsize

\section*{Outputs}

\footnotesize
\begin{longtable}{@{}P{2.0cm}P{3.1cm}P{5.1cm}P{5.5cm}@{}}
\caption{Outputs of the model.}\\
\toprule
\textbf{Symbol} & \textbf{Type} & \textbf{Represents} & \textbf{Obtained how} \\
\midrule
\endfirsthead
\toprule
\textbf{Symbol} & \textbf{Type} & \textbf{Represents} & \textbf{Obtained how} \\
\midrule
\endhead
\bottomrule
\endfoot
$e_i$ & Discrete, $\{0,1\}$. & The challenger's decision to invest (enter). &
  Recorded by the observer as whether the challenger paid $c$, after the
  profile is known and before the draws. \\ \addlinespace
$S$ & Subset of $\{1,\dots,Q\}$. & The entrant set, $S=\{i:e_i=1\}$. &
  Recorded. One play shows the one equilibrium set played, and not the other
  sets the game admits. \\ \addlinespace
$k^*$ & Integer in $\{0,\dots,Q\}$. & The entry count, the largest $k$ with
  the entry condition holding at every rank $j\le k$. By Proposition (count
  invariance) it is the size of every pure-strategy equilibrium entrant
  set. & Recorded as $|S|$. Its definition runs through $\kappa$ and $\Delta$,
  which are constructs, and the recorded count is what that definition
  rationalises. \\ \addlinespace
$w_{(k^*)}$ & Real number. & The marginal entrant, the rank-$k^*$ challenger,
  and her wealth. & Computed from the recorded profile and recorded $k^*$ as
  the wealth of the $k^*$-th richest challenger, whether or not she is in the
  recorded $S$. Under identity multiplicity she can be absent. \\
  \addlinespace
$\{(1),\dots,(k^*)\}$ & Subset of size $k^*$. & The assortative set, the
  $k^*$ richest, equivalently the $k^*$ with the lowest $\kappa$. & Computed
  from the recorded profile and $k^*$, and compared directly with the recorded
  $S$. \\ \addlinespace
$\kappa(w_{(k^*+1)})>\Delta(k^*-1)$ & An inequality. & The condition of
  Proposition (when identities are pinned), under which the assortative set is
  the unique equilibrium. & Not observed, being a statement about constructs.
  In the ratio form $\kappa(w_{(k^*+1)})/V>\Delta(k^*-1)/V$ it survives the
  transformation of $u$ and $V$. \\ \addlinespace
$\sigma_i$ & Real number. & The score, ``what the prize is awarded on.'' &
  Recorded after the draws. \\ \addlinespace
$\mathbf 1\{i\text{ wins}\}$ & Discrete, $\{0,1\}$. & Who takes the prize, the
  holder of the highest score. & Recorded. \\ \addlinespace
$U_i$ & Real number, in the units of $u$. & The payoff,
  $u(w_i-ce_i)+V\cdot\mathbf 1\{i\text{ wins}\}$. The modeller's construct,
  attributed to the agent. & Not observed, being a construct. It rationalises
  the recorded entry decision. Its level is not fixed by the transformation,
  since $b$ is free. \\ \addlinespace
$H_m$; $M_m$ & A CDF; a random variable. & The opponent maximum, the ``CDF of
  the best rival score when $m$ other challengers invest''. $M_m$ is that best
  rival score. & $H_m=F^mG^{Q-1-m}C$ is computed from $F$, $G$, $C$, $Q$, $m$.
  Realisations of $M_m$ are recorded as the best rival score in a contest with
  $m$ other investors. \\ \addlinespace
$\Delta(m)$ & Real number $\ge0$, in the units of $u$; a non-increasing
  function of $m$ (Corollary (monotonicity)). & The gain from investing, the
  right-hand side of the entry condition. A function of the entrant count
  alone (Proposition (anonymity)). & $\Delta(m)$ itself is a construct, since
  it carries $V$. The ratio
  $\Delta(m)/V=\int H_m\,dF-\int H_m\,dG=\mathbb E[\varphi(M_m)]$ (P1) carries
  no utility and is computed from the recorded $F$, $G$, $C$, $Q$. \\
  \addlinespace
$\Delta(m{+}1)-\Delta(m)$; $K$ & Real number $\le0$; $K=F^mG^{Q-2-m}C$, a
  CDF product. & The step of the gain schedule, equal to
  $-\tfrac V2\int\varphi^2\,dK$ (Theorem (key result)). & Computed only, as a
  multiple of $V$, from $F$, $G$, $C$, $Q$. \\ \addlinespace
$\sum_{i\in S}\kappa(w_i)$; $k^*c$ & Real number, in the units of $u$; real
  number, resource units. & Total entry cost in utility units, and its
  resource counterpart, which is identical across equilibria. & The utility
  total is a construct, and summing it across agents is the interpersonal
  comparison the manuscript calls a normative commitment. $k^*c$ is computed
  from the recorded $k^*$ and $c$. \\ \addlinespace
$\sum_iU_i$ & Real number, in the units of $u$. & The unweighted sum of
  payoffs by which Proposition (assortative selection) ranks equilibria. & Not
  observed, being a construct summed across agents by a normative commitment
  (the manuscript's words). The payoffs summed are those that rationalise the
  entry choices. Reading their sum as welfare is a second link, which
  \citet{Sen1973}, pp.~254 and~259, says cannot in general be kept together
  with the first. \\ \addlinespace
$\Delta(m)\to V/(m{+}2)$ as $\mu\to0$ & A function of $m$, independent of
  $Q$. & The P6 benchmark, with the incumbent investing. The manuscript writes
  $1/(m{+}2)$ without the factor $V$. & A computed reference. It is a limit,
  not reached at any recorded $\mu>0$. \\ \addlinespace
$\Delta(m)\le V/(m{+}1)$; $k^*\le V/\kappa$ & Bounds, uniform in $Q$. & P7,
  saturation. & Statements about constructs. In ratio form,
  $\Delta(m)/V\le1/(m{+}1)$ is a computed reference from the recorded $F$,
  $G$, $C$, $Q$, whereas $k^*\kappa/V\le1$ is not observed. \\ \addlinespace
Direction of $k^*$ in $V$ and $c$ & Sign of a weak change. & P8 states that
  $k^*$ is non-decreasing in $V$ and non-increasing in $c$. & Direct
  comparison of recorded $k^*$ across contests that differ only in the
  recorded $c$, or only in the prize, with the prizes ordered by the agents'
  recorded choice between them (the $V$ row of the Inputs table). \\
  \addlinespace
Direction of $k^*$ under a change of profile & Sign of a weak change, rise
  or fall. & Theorem (P9), Theorem (P9, general form), Proposition
  (endogenous margin), Proposition (margin condition). & Direct comparison of
  recorded $k^*$ under two recorded profiles, with $Q$, $c$, the prize and the
  score technology unchanged. The hypothesis ($w_{(k^*)}$ against $x_0$, or
  the margin condition, which reads $d_{k^*}$ for the rise branch and
  $d_{k^*+1}$ for the fall branch) is computed from the two profiles and the
  pre-change $k^*$. No computed reference is needed. \\ \addlinespace
$d_j$; the band & Real numbers; a yes-or-no property of the margin, holding
  when $d_{k^*}<0<d_{k^*+1}$. & $d_j=w'_{(j)}-w_{(j)}$, rank $1$ richest. The
  margin condition is silent exactly inside the band, a pivot-spread never
  reaches the band, and inside the band the displacement signs determine
  nothing, for mean-preserving spreads too (the band). & Computed from two
  recorded sorted profiles and the recorded pre-change $k^*$. Only $d_{k^*}$
  and $d_{k^*+1}$ are needed. \\ \addlinespace
$1-\Lambda(\bar\mu)$ & Real number in $[0,1]$. & The entry rate at which the
  sign of Theorem (P9) flips. & Computed from recorded wealths, and compared
  with the recorded entry rate. \\ \addlinespace
$k^*(\lambda)$; $\bar\lambda$ & A step function of $\lambda$; a real
  number. & Proposition (strictness) states that $k^*$ is non-increasing in
  $\lambda$, and strictly below $k^*(1)$ for $\lambda>\bar\lambda$. &
  $k^*(\lambda)$ is recorded only at the profiles observed. $\bar\lambda$ is
  not observed, since the manuscript asserts its existence only, through
  $\kappa$ and $\Delta(k^*-1)$. \\ \addlinespace
$\Delta(0,Q{+}1)-\Delta(0,Q)$; $W$ & Real number; $W=G^{Q-1}(1-G)F$, a
  function. & P-MU states that the difference equals
  $-V\!\int W\,(dF-dG)$, with the incumbent investing ($C=F$), and that
  $\Delta(0)$ increases in $Q$ iff $\int W\,dF<\int W\,dG$. & A computed
  reference, from the recorded $F$, $G$ and $Q$. It is the change in the gain
  schedule rather than the recorded change in $k^*$ across contests of
  different $Q$. \\ \addlinespace
$\mu^*$ & Real number; distribution-dependent. & The value of $\mu$ at which
  the sign of $d\Delta(0)/dQ$ reverses (the R2 subsection, P-MU). & Computed
  only; never observed. \\
\end{longtable}
\normalsize

\section*{Structure}

\begin{itemize}
\item \emph{Entering only through others.} $u$ enters only through $\kappa$;
  wealth enters only through $\kappa(w_{(j)})$ (Proposition (anonymity));
  $\mu$ and the laws of $r$ and $s$ enter only through $F$ and $G$; $V$
  enters only as a scale factor of $\Delta$, and the entry condition is
  unchanged under $(u,V)\mapsto(au+b,aV)$, $a>0$; $\Lambda$ enters only
  through the realised profile, $\bar\mu$, $1-\Lambda(\bar\mu)$ and quantile
  functions.
\item \emph{Distributions.} $\Lambda$ across challengers, and the laws of $r$
  and $s$, which reach the results only as $F$ and $G$.
\item \emph{Needing a computed reference.} $\Delta/V$ and its step, the P6
  benchmark, the first P7 bound in ratio form, the P-MU sign, $\mu^*$, and the
  flip entry rate $1-\Lambda(\bar\mu)$.
\item \emph{Given directly by recorded comparisons.} $F\le G$ (recorded
  scores of investors against non-investors); independence of the draws from
  wealth (recorded scores across wealth ranks, within one investment status);
  the recorded $S$ against the assortative set; the direction of $k^*$ across
  contests differing only in $c$; the direction of $k^*$ across two recorded
  profiles, with its hypothesis checked on the same records.
\item \emph{Facts that exist but are never observed.} $\mu$, $r_i$, $s_i$,
  and $\Lambda$ as a law.
\item \emph{Constructs, not observed.} $u$, $\kappa$, $V$, $\Delta$ in its
  own units, $U_i$, $\sum_iU_i$, $\sum_{i\in S}\kappa(w_i)$, the condition of
  Proposition (when identities are pinned), the second P7 bound, and
  $\bar\lambda$. Each rationalises recorded entry decisions, or a recorded
  choice between prizes. Under $(u,V)\mapsto(au+b,aV)$ the quantities that do
  not change are exactly the functions of $\kappa/V$ (and of $\Delta/V$, which
  carries no utility). $U_i$ and $\sum_iU_i$ move with $a$ and $b$, whereas
  $\sum_{i\in S}\kappa(w_i)$ moves with $a$ only. The difference of either sum
  between two equilibria keeps its sign. Summing across agents is the
  interpersonal comparison the manuscript calls a normative commitment.
\end{itemize}

\section*{Recovery formulas}

\begin{enumerate}
\item \emph{Linear family $T_\lambda(w)=x_0+\lambda(w-x_0)$.} From two recorded
  sorted profiles, take ranks $a,b$ with $d_a=w'_{(a)}-w_{(a)}$ and
  $d_b=w'_{(b)}-w_{(b)}$. Then
  \[
    \lambda=1+\frac{d_a-d_b}{w_{(a)}-w_{(b)}},\qquad
    x_0=\frac{d_b\,w_{(a)}-d_a\,w_{(b)}}{d_b-d_a}.
  \]
  The formula for $\lambda$ needs $w_{(a)}\ne w_{(b)}$. The formula for $x_0$
  needs $d_a\ne d_b$, equivalently $\lambda\ne1$ and $w_{(a)}\ne w_{(b)}$. Both
  use four recorded wealths.
\item \emph{Linear mean-preserving spread of Theorem (P9).} $\bar\mu$ is the
  mean of the recorded profile, and the spread about $\bar\mu$ preserves that
  mean. Then $\lambda=(w'_{(j)}-\bar\mu)/(w_{(j)}-\bar\mu)$ for any rank $j$
  with $w_{(j)}\ne\bar\mu$. The formula uses one recorded profile and one rank
  of the other.
\item \emph{Gain schedule, up to $V$.} $\Delta(m)/V=\int H_m\,dF-\int
  H_m\,dG=\mathbb E[\varphi(M_m)]$, from the recorded $F$, $G$, $C$ and $Q$
  (P1). The formula needs the regularity of the primitives ($F$, $G$
  continuous, common lower endpoint).
\item \emph{Ratios that rationalise a recorded equilibrium, as a set rather
  than a value.} For a recorded equilibrium $S$ with $|S|=k$,
  $\kappa(w_i)/V\le\Delta(k-1)/V$ for $i\in S$ (when $k\ge1$), and
  $\kappa(w_i)/V>\Delta(k)/V$ for $i\notin S$ (when $k\le Q-1$). These are the
  equilibrium conditions of Proposition (count invariance) divided by $V>0$,
  with $\Delta/V$ from Recovery~3.
\end{enumerate}

\begin{thebibliography}{1}
\bibitem[Schoemaker(1982)]{Schoemaker1982}
Paul J.~H. Schoemaker.
\newblock The expected utility model: its variants, purposes, evidence and
  limitations.
\newblock \emph{Journal of Economic Literature}, 20(2), pp.~529--563, 1982.
\bibitem[Sen(1973)]{Sen1973}
Amartya Sen.
\newblock Behaviour and the concept of preference.
\newblock \emph{Economica}, New Series, 40(159), pp.~241--259, 1973.
\end{thebibliography}

\end{document}
